\chapter*{Abstract}
\addcontentsline{toc}{chapter}{Abstract}

Large language models deployed in serverless or multi-tenant inference environments are frequently evicted from GPU memory when idle and must reload their checkpoints on demand. This loading step can dominate cold-start latency, but the best loading strategy is not obvious: modern storage devices offer high raw throughput, yet practical checkpoint loading performance depends on request structure, checkpoint layout, and software overhead above the storage layer.

This project presents TensorStore, a checkpoint loading framework that supports multiple backends and checkpoint layouts and evaluates them under controlled conditions. The framework supports SafeTensors and a partitioned ServerlessLLM-style layout, with eager-loading backends spanning synchronous POSIX I/O, Tokio-based async loading, and Linux io\_uring. Python bindings and a custom vLLM loader path extend the evaluation beyond Rust-only microbenchmarks.

The final results show that practical checkpoint loading should be framed as an adaptive heuristic problem rather than a one-backend-wins problem. At the Rust layer, smaller eager checkpoints favour synchronous loading, while larger multi-shard checkpoints favour the final multi-worker io\_uring backend. The partitioned ServerlessLLM layout introduces additional nuance, but the dominant practical pattern remains size-sensitive backend selection. Python measurements show that binding-layer overhead is large enough to matter, and that binding-specific optimisation materially improves the public default API. The completed vLLM matrix shows that backend choice affects model initialisation and time-to-first-token far more than steady-state decode. The main conclusion is therefore that LLM checkpoint loading should use adaptive heuristics --- in practice, sync for small models and io\_uring for larger ones --- rather than a fixed backend policy.
